\subsection{DEFINITION OF AN EQUIVALENCE RELATION}

Let \(R\{x, y\}\) be a relation, x and y being distinct letters. The relation R is said to be symmetric (with respect to the letters x and y) if

\[
\mathrm{R} \{x, y \} \Longrightarrow \mathrm{R} \{y, x \}.
\]

If so, by substituting for x and y two letters \(x'\) and \(y'\) , distinct from each other and from all the letters which appear in R, and then by substituting y and x for \(x'\) and \(y'\) , respectively, we see that

\[
\mathrm{R} \left\{y, x \right\} \Longrightarrow \mathrm{R} \left\{x, y \right\}.
\]

Hence \(\mathbb{R}\{x, y\}\) and \(\mathbb{R}\{y, x\}\) are equivalent. Let \(z\) be a letter which does not appear in \(\mathbb{R}\) . The relation \(\mathbb{R}\) is said to be transitive (with respect to the letters \(x, y\) ) if we have

\[
(\mathrm{R} \{x, y \} \text {   and   } \mathrm{R} \{y, z \}) \Longrightarrow \mathrm{R} \{x, z \}.
\]

Examples. The relation x = y is symmetric and transitive. The relation \(X \subset Y\) is transitive but not symmetric. The relation \(X \cap Y = \varnothing\) is symmetric but not transitive.

If \(\mathbf{R}\{x, y\}\) is both symmetric and transitive it is said to be an equivalence relation (with respect to the letters \(x\) and \(y\) ). In this case the notation \(x \equiv y \pmod{\mathbf{R}}\) is sometimes used as a synonym of \(\mathbf{R}\{x, y\}\) ; it is read `` \(x\) is equivalent to \(y\) modulo \(\mathbf{R}\) ''. If \(\mathbf{R}\) is an equivalence relation, we have \(\mathbf{R}\{x, y\} \Longrightarrow (\mathbf{R}\{x, x\} \text{ and } \mathbf{R}\{y, y\})\) , because \(\mathbf{R}\{x, y\}\) implies \(\mathbf{R}\{y, x\}\) , and \((\mathbf{R}\{x, y\} \text{ and } \mathbf{R}\{y, x\})\) implies \((\mathbf{R}\{x, x\} \text{ and } \mathbf{R}\{y, y\})\) by virtue of the definitions.

Let \(\mathbf{R}\{x, y\}\) be a relation; it is said to be reflexive on \(\mathbf{E}\) (with respect to the letters \(x, y\) ) if the relation \(\mathbf{R}\{x, x\}\) is equivalent to \(x \in \mathbf{E}\) . If there is no possible ambiguity about \(\mathbf{E}\) , one says simply, by abuse of language, that \(\mathbf{R}\) is reflexive.

¶ An equivalence relation on E is defined to be an equivalence relation which is reflexive on E. If R\{x, y\} is an equivalence relation on E, we have R\{x, y\} \(\Longrightarrow\) ((x, y) ∈ E × E), hence R has a graph (with respect to the letters x, y). Conversely, suppose that the equivalence relation R\{x, y\} has a graph G. Observe that the relation R\{x, x\} is equivalent to the relation (∃y)R\{x, y\}; for the former implies the latter (Chapter I, §4, no. 2, scheme S5), and conversely, since R\{x, y\} implies R\{x, x\}, (∃y)R\{x, y\} implies (∃y)R\{x, x\} and therefore also R\{x, x\}. Thus R\{x, x\} is equivalent to \(x \in pr_{1}G\) , and hence R is an equivalence relation on \(pr_{1}G\) .

An equivalence on a set E is a correspondence whose source and target are both equal to E, and whose graph F is such that the relation \((x, y) \in \mathbb{F}\) is an equivalence relation on E.

\minorheading{Examples}

\begin{enumerate}
\def\labelenumi{(\arabic{enumi})}
\item
  The relation \(x = y\) is an equivalence relation which has no graph, for if it did, the first projection of this graph would be the set of all objects.
\item
  The relation ``x = y and \(x \in E\) '' is an equivalence relation on E whose graph is the diagonal of \(E \times E\) .
\item
  The relation ``there exists a bijection of X onto Y'' is an equivalence relation which has no graph (cf.~Chapter III, § 3).
\item
  The relation `` \(x \in E\) and \(y \in E\) '' is an equivalence relation on \(E\) , whose graph is \(E \times E\) .
\item
  Suppose \(\mathbf{A} \subset \mathbf{E}\) ; then the relation
\end{enumerate}

\[
(x \in \mathrm{E} - \mathrm{A} \text { and } y = x) \text { or } (x \in \mathrm{A} \text { and } y \in \mathrm{A})
\]

is an equivalence relation on E.

\begin{enumerate}
\def\labelenumi{(\arabic{enumi})}
\setcounter{enumi}{5}
\tightlist
\item
  * The relation `` \(x \in \mathbf{Z}\) and \(y \in \mathbf{Z}\) and \(x - y\) is divisible by 4'' is an equivalence relation on \(\mathbf{Z}_{*}\) .
\end{enumerate}

PROPOSITION 1. A correspondence \(\Gamma\) between X and X is an equivalence on X if and only if it satisfies the following conditions: (a) X is the domain of \(\Gamma\) ; (b) \(\Gamma = \overline{\Gamma}^1\) ; (c) \(\Gamma \circ \Gamma = \Gamma\) .

Let \(\Gamma\) be a correspondence between X and X, and let G be its graph. If \(\Gamma\) is an equivalence on X, then \((x, x) \in G\) for all \(x \in X\) ; hence X is the domain of \(\Gamma\) . The relation \((x, y) \in G\) is equivalent to \((y, x) \in G\) , hence to \((x, y) \in \overline{G}\) , so that \(G = \overline{G}\) and therefore \(\Gamma = \overline{\Gamma}\) . The relations \((x, y) \in G\) and \((y, z) \in G\) imply \((x, z) \in G\) , so that \(G \circ G \subset G\) ; conversely, \((x, y) \in G\) implies \((x, x) \in G\) and therefore \((x, y) \in G \circ G\) , so that \(G \subset G \circ G\) ; hence \(G = G \circ G\) and consequently \(\Gamma = \Gamma \circ \Gamma\) .

Conversely, suppose that conditions (a), (b), and (c) are satisfied. The relation \((x, y) \in G\) is symmetric (by (b)) and transitive (by (c)); hence it is an equivalence relation, and by (a) it is an equivalence relation on \(X\) .

